% Part: counterfactuals
% Chapter: minimal-change-semantics
% Section: true-false

\documentclass[../../../include/open-logic-section]{subfiles}

\begin{document}

\olfileid{cnt}{min}{tf} % BN-SRC-424: প্রতিবাস্তব অংশের সঠিক file-id।

\olsection{প্রতিবাস্তব শর্তবচনের সত্যতা ও মিথ্যাত্ব}

% BN-SRC-770: ছবিগুলির নিকটতম-জগৎ ব্যাখ্যায় ক্ষুদ্রতম পূর্বপক্ষ-স্বীকারক গোলক আছে; সাধারণ সংজ্ঞা তা দাবি করে না।
নিচের ছবিগুলির মতো ক্ষুদ্রতম পূর্বপক্ষ-স্বীকারক গোলক থাকলে
\olref{fig:true}-এ দেখানো মতো নিকটতম $!A$-জগৎগুলির সবই
$!B$-জগৎ হলে প্রতিবাস্তব শর্তবচন $!A \cif !B$
(অশূন্যার্থে) সত্য।
\begin{figure}
\begin{center}
\begin{tikzpicture}[scale=.7]\small
  \spheresystem{5}
  \draw (0,0) node {$w$};
  \propositionintersect{0}{4}{40}{3.5}
  {\tikzset{proposition/.append={smooth,tension=1.4}}
  \proposition[shift={(1.6,-1)}]{90}{1}{30}{4}}
  \path (1.5,3) node[below] {$\formula{B}$};
  \path (3,0) node {$\formula{A}$};
\end{tikzpicture}
\caption{অশূন্যার্থে সত্য প্রতিবাস্তব শর্তবচন}
\ollabel{fig:true}
\end{center}
\end{figure}
আবার~$w$-কে ঘিরে গোলকব্যবস্থায় কোনও $!A$-স্বীকারক গোলকই
না থাকলেও প্রতিবাস্তব শর্তবচনটি $w$-তে সত্য। সে ক্ষেত্রে
এটি \emph{শূন্যার্থে} সত্য (দেখুন \olref{fig:vacuous})।
\begin{figure}
\begin{center}
\begin{tikzpicture}[scale=.7]\small
  \spheresystem{5}
  \draw (0,0) node {$w$};
  \proposition{0}{6.5}{40}{3.5}
  {\tikzset{proposition/.append={smooth,tension=1.4}}
  \proposition[shift={(1.2,-1)}]{90}{1}{30}{4}}
  \path (1.5,3) node[below] {$\formula{B}$};
  \spherepos{0}{7}{node {$\formula{A}$}}
\end{tikzpicture}
\caption{শূন্যার্থে সত্য প্রতিবাস্তব শর্তবচন}
\ollabel{fig:vacuous}
\end{center}
\end{figure}

নিকটতম পূর্বপক্ষ-জগৎ থাকা এই ক্ষেত্রে দুই ভাবে প্রতিবাস্তব
শর্তবচন মিথ্যা হতে পারে। নিকটতম জগৎ না-থাকলে সাধারণ
গোলক-সংজ্ঞাই প্রযোজ্য। এখানে একটি ধরন হল,
নিকটতম $!A$-জগৎগুলির সবগুলি $!B$-জগৎ নয়, কিন্তু কয়েকটি
তা-ই। এ ক্ষেত্রে $!A \cif \lnot !B$-ও মিথ্যা
(দেখুন \olref{fig:false})।
\begin{figure}
\begin{center}
\begin{tikzpicture}[scale=.7]\small
  \spheresystem{5}
  \draw (0,0) node {$w$};
  \propositionintersect{0}{4}{40}{3.5}
  {\tikzset{proposition/.append={smooth,tension=1.4}}
  \proposition[shift={(1.6,0)}]{90}{1}{30}{3}}
  \path (1.5,3) node[below] {$\formula{B}$};
  \path (3,0) node {$\formula{A}$};
\end{tikzpicture}
\caption{মিথ্যা প্রতিবাস্তব শর্তবচন, বিপরীতটিও মিথ্যা}
\ollabel{fig:false}
\end{center}
\end{figure}
নিকটতম $!A$-জগৎগুলির সঙ্গে $!B$-জগৎগুলির কোনও মিল না
থাকলে $!A \cif !B$ মিথ্যা। কিন্তু সে ক্ষেত্রে নিকটতম সব
$!A$-জগৎই $\lnot !B$-জগৎ; ফলে $!A \cif \lnot !B$ সত্য
(দেখুন \olref{fig:false-opposite})।
\begin{figure}
\begin{center}
\begin{tikzpicture}[scale=.7]\small
  \spheresystem{5}
  \draw (0,0) node {$w$};
  \propositionintersect{0}{4}{40}{3.5}
  {\tikzset{proposition/.append={smooth,tension=1.4}}
  \proposition[shift={(-1,0)}]{90}{1}{30}{3}}
  \path (-1,3) node[below] {$\formula{B}$};
  \path (3,0) node {$\formula{A}$};
\end{tikzpicture}
\caption{মিথ্যা প্রতিবাস্তব শর্তবচন, বিপরীতটি সত্য}
\ollabel{fig:false-opposite}
\end{center}
\end{figure}

আগের বিভাগের S5-তে কঠোর শর্তবচনের বিপরীতে প্রতিবাস্তব
শর্তবচন আপতিক হতে পারে।
\olref{fig:contingent}-এর গোলক মডেলটি দেখুন। $u$-এর নিকটতম
$!A$-জগৎগুলির সবই $!B$-জগৎ; তাই $\mSat{M}{!A \cif
  !B}[u]$। কিন্তু $v$-এর নিকটতম কিছু $!A$-জগৎ
$!B$-জগৎ নয়; তাই $\mSat/{M}{!A \cif !B}[v]$।
\begin{figure}
\begin{center}
\begin{tikzpicture}[scale=.6]\small
  \spheresystem{7}
  \spheresystem[shift={(0:3.1)},dashed]{4}
  \propositionintersect{45}{4}{40}{4.5}
  \begin{scope}
    \clip \propositionplot{45}{4}{40}{4.5} ;
    \spherelayer[shift={(0:3.1)}]{3}
  \end{scope}
  \proposition[shift={(1.4,-.5)}]{90}{1}{35}{5}
  \path (0,0) node {$u$};
  \path (0:3.1) node {$v$};
  \spherepos{35}{9}{node {$\formula{A}$}}
  \spherepos{80}{9}{node {$\formula{B}$}}
\end{tikzpicture}
\end{center}
\caption{আপতিক প্রতিবাস্তব শর্তবচন}
\ollabel{fig:contingent}
\end{figure}
\end{document}
